\appendix

\section[Monotonicities along the p-capacitary potential]{Monotonicities along the $\boldsymbol{p}$-capacitary potential}
\label{app:monotonicities}
Here we slightly improve the monotonicity results in~\cite{agostiniani_riemannianpenroseinequalitynonlinear_2022,chan_monotonicitygreenfunctions_2022}. Inspired by these two works, we are approximating the $p$-capacitary potential with a family of smooth functions. To enter more in detail, let $(M,g)$ be a strongly $p$-nonparabolic Riemannian manifold with (possibly empty) boundary. Let $\Omega \subset M$ be homologous to $\partial M$ and $u_p$ is the solution to~\eqref{eq:pb-intro} starting at $\Omega$. For every $T>1$ let $\Omega_T$ be strictly homologous to $\partial M$ with connected boundary and containing $\set{u_p>\alpha_p(T)}$, where $\alpha_p(T) = T^{-(3-p)/(p-1)}$. Then, we define $u_p^\varepsilon$ as the solution to the following boundary value problem:
\begin{equation}\label{eq:pb-intro-eps}
 \begin{cases}
 \Delta_{p}^\varepsilon u^\varepsilon_p = 0 &\text{on $\Int\Omega_T \smallsetminus \Omega $,}\\
 u_p^\varepsilon=1 & \text{on $\partial \Omega$,}\\
 u_p^\varepsilon=u_p & \text{on $\partial\Omega_T$,}
 \end{cases}
\end{equation}
where
\[
 \Delta_{p}^\varepsilon f = \div\big( \abs{ \D f}_\varepsilon^{p-2} \D f\big) \qquad \text{and} \qquad {\abs{}}_\varepsilon= \sqrt{ \abs{}^2 +\varepsilon^2}.
\]
The function $u^\varepsilon_p$ is smooth away from the exterior boundary and converges in $\CS^{1,\beta}_{\loc}$ to the $p$-capacitary potential $u_p$ as $\varepsilon \to 0^+$. Indeed, this family was used in~\cite{dibenedetto_alphalocalregularityweak_1983} to prove $\CS^{1,\beta}_{\loc}$-regularity of $p$-harmonic functions. Moreover, looking more carefully at the proof of~\cite[Lemma~2.1]{lou_singularsetslocalsolutions_2008}, $\abs{\D u^\varepsilon_p}^{p-1}$~is uniformly bounded in $W^{1,2}_{\loc}$. Hence, up to a not relabeled subsequence, we can always assume that $\abs{ \D u_p^\varepsilon}^{p-1}$ weakly converges in $W^{1,2}_{\loc}$. Moreover, since $\abs{ \D u_p^\varepsilon}$ converges uniformly to~$| \D u_p |$, the weak limit of $\D \abs{ \D u_p^\varepsilon}^{p-1}$ must be $\D | \D u_p |^{p-1}$.

We are now ready to prove Theorem~\ref{thm:AMMO_monotonicity}.

\begin{proof}[Proof of Theorem~\ref{thm:AMMO_monotonicity}] For ease of computations, we rewrite the function $t\mapsto \ma_{H}^{\tp}(\partial \Omega_t)$ in terms of the $p$-capacitary potential, that is,
\[
 U_p(t)= 4 \pi t + \frac{(p-1)^2}{(3-p)^2} t^{\frac{5-p}{p-1}} \int_{\set{u_p = \alpha_p(t)}} | \D u_p |^2 \dif \sigma -\frac{(p-1)}{(3-p)} t^{\frac{2}{p-1}} \int_{\set{u_p = \alpha_p(t)}} | \D u_p |\H \dif \sigma.
\]
Observe that
\[
 \frac{\ncapa_p(\partial \Omega)^\frac{1}{3-p}}{8\pi} U_p\big(\ee^{\frac{t}{3-p}}\big) = \ma_{H}^{\tp}(\partial \Omega_t)= \frac{1}{8\pi}F_p\big(\ncapa_p(\partial \Omega)^\frac{1}{3-p} \ee^{\frac{t}{3-p}}\big),
\]
where $F_p$ is the function defined in~\cite{agostiniani_riemannianpenroseinequalitynonlinear_2022}. Equation~\eqref{eq:AMMO_monotonicity} now follows from computations in~\cite[Section 1.2]{agostiniani_riemannianpenroseinequalitynonlinear_2022}. The function $U_p$ is well defined on $[0,+\infty)$. Indeed, one can observe that
\[
 \abs{ \H} | \D u_p |^2 = | \D u_p |^{3-p} | \D | \D u_p |^{p-1}| \in L^2_{\loc}(M\smallsetminus \Omega),
\]
since $1<p<3$ and $| \D u_p |^{p-1} \in L^\infty_{\loc}\cap W^{1,2}_{\loc} (M\smallsetminus \Omega)$ by~\cite[Lemma~2.1]{lou_singularsetslocalsolutions_2008}. Then, by coarea formula, the function
\[
 t \mapsto \int_{\set{u_p= \alpha_p(t)}} | \D u_p | \H \dif \sigma\in L^1_{\loc}(0,+\infty)
\]
and its equivalence class does not depend on the representative of $| \D u_p | \H $. In particular, the function $t\mapsto \ma_{H}^{\tp}(\partial \Omega_t) \in L^1_{\loc}(0,+\infty)$.

It only remains to prove that it has nonnegative first derivative in the sense of distributions, which both gives that $t\mapsto \ma_{H}^{\tp}(\partial \Omega_t) \in \BV_{\loc}(0,+\infty)$ and that admits a nondecreasing representative. Fix $T>0$ and $u^\varepsilon_p$ be the solution to~\eqref{eq:pb-intro-eps}. One can now define the function
\[
 F_p^\varepsilon(t)= 4 \pi t + \frac{(p-1)^2}{(3-p)^2} t^{\frac{5-p}{p-1}} \int_{\set{u^\varepsilon_p = \alpha_p(t)}} \abs{ \D u^\varepsilon_p}^2 \dif \sigma -\frac{(p-1)}{(3-p)} t^{\frac{2}{p-1}} \int_{\set{u^\varepsilon_p = \alpha_p(t)}} \abs{ \D u^\varepsilon_p}\H^\varepsilon\dif \sigma
\]
for every $t \in (1,T)$, where $\H^\varepsilon$ is the mean curvature of the level $\set{u^\varepsilon_p = \alpha_p(t)}$. By~\cite[Lemma~1.2]{agostiniani_riemannianpenroseinequalitynonlinear_2022}, we have that the function $F^\varepsilon_p$ is almost monotone, in the sense that
\[
 F^\varepsilon_p(t)- F^\varepsilon_p(s) \geq -\varepsilon
 \frac{(p+1)^2(p-1)}{(3-p)^3} \int_{ \set{\alpha_p(t) < u^\varepsilon_p<\alpha_p(s)}}\abs{ \D u_p^\varepsilon}^2(u^\varepsilon_p)^{-\frac{p-1}{3-p}-3} \dif \mu
\]
holds for almost every $0\leq s\leq t \leq T$.
In particular, by coarea formula and Lebesgue differentiation theorem,
\[
 (F^\varepsilon_p)'(t) \geq
 - \varepsilon \frac{(p+1)^2}{(3-p)^2}t^{\frac{2(3-p)}{p-1}} \int_{\set{u^\varepsilon_p=\alpha_p(t)}} \abs{ \D u^\varepsilon_p} \dif \sigma
\]
holds for almost every $t \in (1,T)$.

The monotonicity follows if one can prove the following claim.
\begin{Claim}
$F_p^\varepsilon$ converges to $F_p$ in the sense of distributions as $\varepsilon\to 0^+$.
\end{Claim}
Indeed, if that is the case, for every nonnegative test function $\varphi \in \CS^\infty_c(1,T)$, we obtain that
\begin{align*}
 -\int_1^T \varphi'(t) F_p(t)\dif t & =-\lim_{\varepsilon\to 0} \int_1^T\varphi'(t) F^p_\varepsilon(t)\dif t \\
 & \geq - \frac{(p+1)^2}{(3-p)^2} \lim_{\varepsilon \to 0} \varepsilon \int_1^T \varphi(t) t^{\frac{2(3-p)}{p-1}} \int_{\set{u^\varepsilon_p=\alpha_p(t)}} \abs{ \D u^\varepsilon_p} \dif \sigma\dif t \\
 & =-\frac{(p+1)^2}{(3-p)^2} \lim_{\varepsilon \to 0} \varepsilon \int_{\Omega_T \smallsetminus \Omega} \varphi \big(\alpha_p^{-1}\big(u^\varepsilon_p\big)\big) \frac{ \big| \D u^\varepsilon_p \big|^2}{\big(u^\varepsilon_p\big)^{\frac{3-p}{p-1}+3}} \dif \mu=0,
\end{align*}
since $u_p^\varepsilon$ converges to $u_p$ in $\CS^{1,\beta}_{\loc}$. This shows that $F_p$ has nonnegative first derivative in the sense of distributions, proving its monotonicity.

We now turn to prove the claim. Consider any $\varphi \in \CS^\infty_c(0,+\infty)$. The first term is independent of $\varepsilon$. As far as the second term is concerned, by coarea formula, we have that
\[
 \int_1^T\varphi(t)t^{\frac{5-p}{p-1}} \int_{\set{u_p^\varepsilon=\alpha_p(t)}} \abs{\D u^\varepsilon_p}^2 \dif \sigma \dif t = \frac{(p-1)}{(3-p)} \int_{\Omega_T \smallsetminus \Omega} \big(u^\varepsilon_p\big)^{-\frac{7-p}{3-p}}\varphi \big(\alpha_p^{-1}\big(u^\varepsilon_p\big)\big) \abs{ \D u_p^\varepsilon}^3 \dif \mu.
\]
Since the function $\varphi$ is smooth with compact support in $(1,T)$ and $u_p^\varepsilon$ converges to $u_p$ in $\CS^{1,\beta}_{\loc}$, the right-hand side converges to
\[
 \frac{(p-1)}{(3-p)} \int_{\Omega_T \smallsetminus \Omega} u_p^{-\frac{7-p}{3-p}}\varphi \big(\alpha_p^{-1}\big(u^\varepsilon_p\big)\big)| \D u_p |^3 \dif \mu = \int_1^T \varphi(t)t^{\frac{5-p}{p-1}} \int_{\set{u_p^\varepsilon=\alpha_p(t)}} |\D u_p |^2 \dif \sigma \dif t,
\]
where the identity follows by coarea formula. The last term is a little trickier since it involves second derivatives of the function $u_p^\varepsilon$ that are not converging uniformly as $\varepsilon \to 0$ to the corresponding ones for $u_p$. Employing again the coarea formula and by straightforward computations, we then have that
\begin{gather*}
 \int_1^T\varphi(t)t^{\frac{2}{p-1}} \int_{\set{u_p^\varepsilon=\alpha_p(t)}} \abs{\D u^\varepsilon_p} \H_\varepsilon\dif \sigma \dif t \\
 \qquad {}=
 \int_1^T\varphi(t)t^{\frac{2}{p-1}} \int_{\set{u_p^\varepsilon=\alpha_p(t)}} \frac{\big\langle \D \big| \D u^\varepsilon_p\big|^{p-1}\big|\D u^\varepsilon_p\big\rangle}{\big| \D u_p^\varepsilon\big|^{p-1} }\Bigg( 1 - \frac{(p-2)}{(p-1)} \frac{ \varepsilon^2}{\big|\D u^\varepsilon_p\big|_\varepsilon^2}\Bigg)\dif \sigma \dif t\\
 \qquad {}=\frac{(p-1)}{(3-p)} \int_{\Omega_T \smallsetminus \Omega} \big(u^\varepsilon_p\big)^{-\frac{4}{3-p}}\varphi \big(\alpha_p^{-1}\big(u^\varepsilon_p\big)\big)\frac{\big\langle\D \big| \D u^\varepsilon_p\big|^{p-1}\big|\D u^\varepsilon_p\big\rangle}{\big| \D u_p^\varepsilon\big|^{p-2} }\Bigg( 1 - \frac{(p-2)}{(p-1)} \frac{ \varepsilon^2}{\big|\D u^\varepsilon_p\big|_\varepsilon^2}\Bigg) \dif \mu.
\end{gather*}
As before we want to prove that the right-hand side converges to the corresponding term for $u_p$. Since $u^\varepsilon_p\to u_p$ in $\CS^{1, \beta}_{\loc}$ and $\D\big| \D u^\varepsilon_p\big|^{p-1}\rightharpoonup \D| \D u_p |^{p-1} $ weakly in $L^2_{\loc}$, we have that
\begin{gather*}
 \lim_{\varepsilon \to 0}\frac{(p-1)}{(3-p)} \int_{\Omega_T \smallsetminus \Omega} \big(u^\varepsilon_p\big)^{-\frac{4}{3-p}}\varphi \big(\alpha_p\big(u^\varepsilon_p\big)\big) \frac{\big\langle\D \big| \D u^\varepsilon_p\big|^{p-1}\big|\D u^\varepsilon_p\big\rangle}{\big| \D u_p^\varepsilon\big|^{p-2} }\dif \mu\\
\qquad {} =\frac{(p-1)}{(3-p)} \int_{\Omega_T \smallsetminus \Omega} (u_p)^{-\frac{4}{3-p}}\varphi \big(\alpha_p^{-1}\big(u^\varepsilon_p\big)\big)\frac{\big\langle\D | \D u_p |^{p-1}|\D u_p\big\rangle}{| \D u_p |^{p-2} }\dif \mu\\
\qquad {} =\int_1^T\varphi(t)t^{\frac{2}{p-1}}\int_{\set{u_p=\alpha_p(t)}} |\D u_p | \H \dif \sigma \dif t.
\end{gather*}
Moreover, the remaining term vanishes. H\"older's inequality and equi-boundedness in $L_{\loc}^2(\Omega_T \smallsetminus \Omega)$ of $\big|\D \big|\D u^\varepsilon_p\big|^{p-1}\big|$ yield
\begin{equation}\label{eq:main_inequality_monotonicity_AMMO}
 \int_{K}\abs{\D u^\varepsilon_p}^{3-p}\big|\D \abs{ \D u^\varepsilon_p}^{p-1}\big|\frac{ \varepsilon^2}{\big| \D u^\varepsilon_p\big|_\varepsilon^2}\dif \mu\leq \kst_1\Bigg(\int_K\frac{\varepsilon^4}{\big| \D u^\varepsilon_p\big|_\varepsilon^4} \abs{\D u^\varepsilon_p}^{6-2p}\dif \mu\Bigg)^\frac{1}{2}
\end{equation}
for every $K$ compactly contained in $\set{ \alpha_p(T) < u_p^\varepsilon <1}$ and for some positive constant $\kst_1$. Observe that
\[
 \abs{\D u^\varepsilon_p}^{6-2p}\frac{\varepsilon^4}{\big| \D u^\varepsilon_p\big|_\varepsilon^4} \leq \abs{\D u^\varepsilon_p}^{6-2p}\leq \kst_2,
\]
since the function $\big| \D u_\varepsilon^p\big|$ converges locally uniformly and $1<p<3$. The left-hand side converges almost everywhere to $0$. Indeed, if a point belongs to critical set of $u_p$, \smash{$\big|\D u_\varepsilon^p\big|^{6-2p}\to 0$} as $\varepsilon\to 0$. Otherwise, $\big| \D u_\varepsilon^p\big|$ is definitely bounded away from $0$, then \smash{$\big| \D u^\varepsilon_p\big|_\varepsilon^4$} is not vanishing, thus the left-hand side is controlled by $\varepsilon^4$ up to a constant. By dominated convergence theorem, the right-hand side in~\eqref{eq:main_inequality_monotonicity_AMMO} approaches $0$ as $\varepsilon \to 0$, so does the left-hand side, concluding the step. \end{proof}

We use this theorem to study the quantity introduced in~\cite[Theorem 1.2]{chan_monotonicitygreenfunctions_2022} along the level set of the solution $w_p$ to~\eqref{eq:pb-intro}.

\begin{Theorem}\label{thm:CCLT_monotonicity}
 Let $(M,g)$ be a strongly $p$-nonparabolic Riemannian $3$-manifold with smooth, compact and connected possibly empty boundary $\partial M$. Assume that $H_2(M, \partial M;\Z) = \set{0}$. Let $\Omega\subseteq M$ be bounded closed with connected $\CS^1$-boundary homologous to $\partial M$ and with $\h\in L^2(\partial \Omega)$. Let $w_p$ be the solution to~\eqref{eq:pb-intro} starting at $\Omega$. Then, denoting $\Omega_t = \set{w_p \leq t}$, the function
 \begin{equation}\label{eq:monotoncity_chan}
 t \mapsto \frac{\ncapa_p(\partial \Omega_t)^{-\frac{1}{p-1}}}{4\pi(3-p)}\Bigg( 4 \pi - \int_{\partial \Omega_t} \frac{\abs{\D w_p}^2}{(3-p)^2}\dif \mu\Bigg)
 \end{equation}
 belongs to $W^{1,1}_{\loc}(0,+\infty)$ and
 \begin{equation}\label{eq:weak_derivative}
 \frac{\dif}{\dif t}\Bigg[\ncapa_p(\partial \Omega_t)^{-\frac{1}{p-1}}\Bigg( 4 \pi - \int_{\partial \Omega_t} \frac{\abs{\D w_p}^2}{(3-p)^2}\dif \sigma\Bigg)\Bigg] =- \frac{8\pi}{p-1}\ncapa_p(\partial \Omega_t)^{-\frac{2}{(3-p)(p-1)}}\ma_{H}^{\tp}(\partial \Omega_t),
 \end{equation}
 for almost every $t \in [0,+\infty)$.
\end{Theorem}

\begin{Remark}\label{rmk:cclt_mon}
 Observe that~\eqref{eq:monotoncity_chan} is, up to multiplying by a constant and changing variables, the quantity studied in~\cite{chan_monotonicitygreenfunctions_2022,munteanu_comparisontheorems3dmanifolds_2023} for $p\in(1,2]$ along the level set of the $p$-Green function. In particular, coupled with Theorem~\ref{thm:AMMO_monotonicity}, the above result implies that the monotonicity property of~\eqref{eq:monotoncity_chan} is preserved if, in place of the $p$-Green's function, the $p$-capacitary potential of a connected $\partial \Omega$ with nonnegative $p$-Hawking mass is considered. Clearly, the monotonicity results in~\cite{chan_monotonicitygreenfunctions_2022, munteanu_comparisontheorems3dmanifolds_2023} are recovered applying Theorem~\ref{thm:CCLT_monotonicity} to the $p$-Green's function. Hence, we settle the question raised in~\cite{chan_monotonicitygreenfunctions_2022} about the monotonicity of their quantities in the range $p \in (2, 3)$.

 Finally, observe that $\tilde{\ma}^{\tp}_H$ is obtained multiplying the function~\eqref{eq:monotoncity_chan} by the $p$-capacity term $\ncapa_p(\partial \Omega_t)^{2/(3-p)(p-1)}$ which is exponentially growing as $t \to +\infty$. This term forces the quantity~$\tilde{\ma}^{\tp}_H$ to be monotone \emph{nondecreasing} under assumption~\eqref{eq:energy_hp} (see Lemma~\ref{thm:inequality_modifiedgeroch_AMMO}) even when~\eqref{eq:monotoncity_chan} is monotone \emph{nonincreasing}.
\end{Remark}

\begin{proof}[Proof of Theorem~\ref{thm:CCLT_monotonicity}] Fix $T>1$ and let $u^\varepsilon_p$ the solution to the problem~\eqref{eq:pb-intro-eps}. Consider any $\varphi \in \CS^\infty_c(0,T)$. Employing coarea formula and integration by parts, we have that
\begin{align}
 \int_1^{T} \varphi'(t) \int_{\set{u^\varepsilon_p=\alpha_p(t)}} \abs{ \D u^\varepsilon_p}^2 \dif \sigma \dif t &{}= \frac{(p-1)}{(3-p)} \int_{\Omega_T \smallsetminus \Omega} \varphi' \big(\big(u_p^\varepsilon\big)^{-\frac{p-1}{3-p}} \big)\big(u_p^\varepsilon\big)^{-\frac{2}{3-p}} \abs{ \D u^\varepsilon_p}^3 \dif \mu \nonumber\\
 &{}=- \int_{\Omega_T \smallsetminus \Omega} \big\langle\big| \D u^\varepsilon_p\big|\D u^\varepsilon_p\big| \D \big[\varphi \big(\big(u_p^\varepsilon\big)^{-\frac{p-1}{3-p}}\big)\big]\big\rangle \dif \mu \nonumber\\
 &{}= \int_{\Omega_T \smallsetminus \Omega} \div \big( \abs{ \D u^\varepsilon_p} \D u^\varepsilon_p\big) \varphi \big(\big(u_p^\varepsilon\big)^{-\frac{p-1}{3-p}}\big) \dif \mu. \label{eq:main_chan}
\end{align}
Clearly, employing $\CS^{1, \beta}_{\loc}$ convergence of $u^\varepsilon_p \to u_p$ as $\varepsilon \to 0$, we obtain that
\[
 \lim_{\varepsilon\to 0}\int_1^{T} \varphi'(t) \int_{\set{u^\varepsilon_p=\alpha_p(t)}} \abs{ \D u^\varepsilon_p}^2 \dif \sigma \dif t = \int_1^{T} \varphi'(t) \int_{\set{u_p=\alpha_p(t)}} | \D u_p |^2 \dif \sigma \dif t.
\]
Moreover, a straightforward computation leads to
\[
 \div \big( \abs{ \D u^\varepsilon_p} \D u^\varepsilon_p\big) =\frac{\big| \D u^\varepsilon_p\big|^{2-p}}{(p-1)} \Bigg( (3-p) \frac{ \big| \D u^\varepsilon_p\big|^2}{\big| \D u^\varepsilon_p \big|^2_\varepsilon}- \frac{ \varepsilon^2}{\big| \D u^\varepsilon_p\big|_{\varepsilon}^2} \Bigg) \big\langle\D \big| \D u^\varepsilon_p \big|^{p-1} \big| \D u^\varepsilon_p\big\rangle.
\]
Arguing as in the previous theorem, since $\big|\D u^\varepsilon_p\big| \to |\D u_p |$ locally uniformly and $\D\big|\D u^\varepsilon_p\big|^{p-1} \rightharpoonup\D |\D u_p |^{p-1}$ weakly $L^2_{\loc}$ as $\varepsilon\to 0$, one gets
\begin{gather*}
 \lim_{\varepsilon \to 0} \int_{\Omega_T \smallsetminus \Omega} \frac{\big| \D u^\varepsilon_p\big|^{2-p}}{(p-1)} \Bigg( (3-p) \frac{ \big| \D u^\varepsilon_p\big|^2}{\big| \D u^\varepsilon_p \big|^2_\varepsilon}- \frac{ \varepsilon^2}{\big| \D u^\varepsilon_p\big|_{\varepsilon}^2} \Bigg) \big\langle\D \big| \D u^\varepsilon_p \big|^{p-1} \big| \D u^\varepsilon_p\big\rangle \varphi \big(\big(u_p^\varepsilon\big)^{-\frac{p-1}{3-p}}\big) \dif \mu\\
 \qquad {}=
 \frac{(3-p)}{(p-1)} \int_{\Omega_T \smallsetminus \Omega} \frac{ \big\langle \D | \D u_p |^{p-1} | \D u_p\big\rangle}{| \D u_p |^{p-2}} \varphi \big((u_p)^{-\frac{p-1}{3-p}}\big)\dif \mu \\
 \qquad {}=
 \frac{(3-p)^2}{(p-1)^2}\int_1^{T}\varphi(t) t^{-\frac{2}{p-1}} \int_{\Omega_T \smallsetminus \Omega} \H \abs{ \D u} \dif \sigma \dif t,
\end{gather*}
where the last equality follows by coarea formula and $\H= \big\langle\D | \D u_p |^{p-1} | \D u_p\big\rangle / | \D u_p |^p$. Then, passing to the limit as $\varepsilon \to 0$ in~\eqref{eq:main_chan}, we obtain
\[
 \int_1^{T} \varphi'(t) \int_{\set{u_p=\alpha_p(t)}} | \D u_p |^2 \dif \sigma \dif t =
 \frac{(3-p)^2}{(p-1)^2}\int_1^{T}\varphi(t) t^{-\frac{2}{p-1} } \int_{\set{u_p= \alpha_p(t)}} \H \abs{ \D u} \dif \sigma \dif t.
\]
By arbitrariness of $T$ and $\varphi$ one has that the function
\[
 t \mapsto H_p(t)=\Bigg(4 \pi t^{-\frac{3-p}{p-1}}-\frac{(p-1)^2}{(3-p)^2}t^{\frac{3-p}{p-1}} \int_{\set{u=\alpha_p(t)}} |\D u_p |^2 \dif \sigma\Bigg)
\]
belongs to $W^{1,1}_{\loc}(1,+\infty)$ and its derivative is given by
\begin{align*}
H'_p(t)={}& -\frac{(3-p)}{(p-1)}t^{-\frac{p+1}{p-1}}\Bigg(4 \pi t+ \frac{(p-1)^2}{(3-p)^2} t^{\frac{5-p}{p-1}} \\
&{}\times \int_{\set{u=\alpha_p(t)}} |\D u_p |^2 \dif \sigma - \frac{(p-1)}{(3-p)} t^{\frac{2}{p-1}}\int_{\set{u_p= \alpha_p(t)}} \H \abs{ \D u} \dif \sigma\Bigg)
\end{align*}
holds for almost every $t \in (1,+\infty)$. Then
\begin{align*}
 \frac{\dif}{\dif t}\Bigg[\ncapa_p(\partial \Omega_t)^{-\frac{1}{p-1}}\Bigg( 4 \pi - \int_{\partial \Omega_t} \frac{\big|\D w_p\big|^2}{(3-p)^2}\dif \mu\Bigg)\Bigg] &{}= \frac{\ncapa_p(\partial \Omega)^{-\frac{1}{p-1}}}{3-p} H'_p\big(\ee^{\frac{t}{3-p}}\big) \ee^{\frac{t}{3-p}}\\
 &{}=- \frac{8\pi}{p-1}\ncapa_p(\partial \Omega_t)^{-\frac{2}{(3-p)(p-1)}}\ma_{H}^{\tp}(\partial \Omega_t),
\end{align*}
concluding the proof.\end{proof}
